\documentclass[DIN, pagenumber=false, fontsize=11pt, parskip=half]{scrartcl}

\usepackage[utf8]{inputenc}
\usepackage[T1]{fontenc}
\usepackage{template}
\usepackage{amsmath}

\DeclareMathOperator{\lcm}{lcm}

\titlehead{\centering\small
    Worked solutions to selected exercises from\\
    Richard Hammack, \emph{Book of Proof, edition 2.2}\\
\rule{0.4\linewidth}{0.4pt}}

\title{}
\author{} % Oscar A. Carrera
\date{}

\begin{document}
\maketitle

\setcounter{section}{8}

\section{Disproof}

\textbf{Each of the following statements is either true or false.
    If a statement is true, prove it. If a statement is false, disprove it.
These exercises are cumulative, covering all topics addressed in Chapters 1-9.}

\begin{problems}
    \problem For every natural number $n$, the integer $2n^2-4n+31$ is prime.
    \begin{disproof}
        The statement is false. 
        For a counterexample, let $n = 31$.
        Then $2n^2-4n+31 = 2(31^2)-4(31)+31 = 1829 = 31 \cdot 59$,
        which is not prime.
    \end{disproof}

    \problem For every natural number $n$, the integer $n^2 + 17n+17$ is prime.
    \begin{disproof}
        The statement is false.
        For a counterexample, let $n = 1$.
        Then $n^2+17n+17 = 1^2 + 17(1) + 17 = 35 = 7 \cdot 5$,
        which is not prime.
    \end{disproof}

    \problem If $A,B,C$ and $D$ are sets, then $(A \times B) \cap (C \times D) = (A \cap C) \times (B \cap D)$.
    \begin{proof}
        Observe the following sequence of equalities.
        \begin{align*}
            \MoveEqLeft (A \times B) \cap (C \times D)\\
            & = \set{(x,y) : (x,y) \in A \times B \land (x,y) \in C \times D} && \text{(definition of $\cap$)}\\
            & = \set{(x,y) : (x \in A \land y \in B) \land (x \in C \land y \in D)} && \text{(definition of $\times$)}\\
            & = \set{(x,y) : (x \in A \land x \in C) \land (y \in B \land y \in D)} && \text{(regrouping)}\\
            & = \set{(x,y) : x \in A \cap C \land y \in B \cap D} && \text{(definition of $\cap$)}\\
            & = \set{(x,y) : (x,y) \in (A \cap C) \times (B \cap D)} && \text{(definition of $\times$)}\\
            & = (A \cap C) \times (B \cap D) && \text{(definition of set-builder)}
        \end{align*}
    \end{proof}

    \problem If $A,B$ and $C$ are sets, then $A - (B \cup C) = (A - B) \cup (A - C)$.
    \begin{disproof}
        The statement is false.
        For a counterexample, let $A = \set{1,2}$, $B = \set{1}$ and $C = \set{2}$.
        Then $A - (B \cup C) = \set{1,2} - \set{1,2} = \emptyset$,
        but $(A - B) \cup (A - C) = \set{2} \cup \set{1} = \set{1,2}$.
        Since $\emptyset \ne \set{1,2}$, the two sides are not equal.
    \end{disproof}

    \problem If $A$ and $B$ are sets and $A \cap B = \emptyset$, then $\mathcal{P}(A) -\mathcal{P}(B) \subseteq \mathcal{P}(A - B)$.
    \begin{proof}
        Suppose $A \cap B = \emptyset$.
        Let $X \in \mathcal{P}(A) - \mathcal{P}(B)$.
        Then $X \in \mathcal{P}(A)$ and $X \notin \mathcal{P}(B)$, so $X \subseteq A$.
        Now we claim $A - B = A$. 
        Let $a \in A - B$.
        Then $a \in A$ and $a \notin B$.
        Thus every element of $A - B$ is an element of $A$, so $A - B \subseteq A$.
        Conversely, let $a \in A$.
        If $a \in B$, then $a \in A$ and $a \in B$, so $a \in A \cap B = \emptyset$, which is impossible.
        Thus $a \notin B$.
        So we have $a \in A$ and $a \notin B$, which means $a \in A - B$.
        Hence $A \subseteq A - B$.
        Therefore $A - B = A$.
        So we have $X \subseteq A = A - B$, thus $X \in \mathcal{P}(A-B)$.
        Since $X$ was arbitrary, $\mathcal{P}(A) - \mathcal{P}(B) \subseteq \mathcal{P}(A-B)$.
    \end{proof}

    \problem If $a,b,c \in \mathbb{N}$ and $ab, bc$ and $ac$ all have the same parity, the $a,b$ and $c$ all have the same parity.
    \begin{disproof}
        The statement is false.
        For a counterexample, consider $a = 2n$, $b = 2m+1$, and $c = 2l$ for some $n,m,l \in \mathbb{N}$. Then $ab = 2(2mn+n)$, $bc = 2(2lm+l)$ and $ac = 2(2nl)$, this means $ab,bc$ and $ac$ are even, so they all have the same parity.
        But $a$ is even, $b$ is odd, and $c$ is even. That is $a$, $b$, and $c$ all do not have the same parity. 
    \end{disproof}

\end{problems}

\end{document}
